\section{Bayesian Compound Identification Confidence
Scoring}\label{bayesian-compound-identification-confidence-scoring}

\textbf{Dataset:} LC-BinBase, HILIC negative-ion TTOF (Premier)
\textbf{Confirmed spectra evaluated:} 1,631 (manually annotated,
non-yy/zz) \textbf{Unconfirmed spectra expanded:} 5,899 (never
previously queried) \textbf{Total library hits scored:} 195,367
candidate matches

\begin{center}\rule{0.5\linewidth}{0.5pt}\end{center}

\subsection{Part I --- The Problem and
Approach}\label{part-i-the-problem-and-approach}

\subsubsection{1.1 What We Are
Solvingxssxs}\label{what-we-are-solvingxssxs}

Standard LC-MS/MS library search returns a ranked list of candidate
compounds with no uncertainty estimate. The top hit might be correct
with 99\% confidence, or it might be one of ten equally plausible
structural isomers with no way to tell them apart. Analysts either
accept all top hits blindly or spend large amounts of time manually
reviewing individual spectra.

This project builds a \textbf{Bayesian evidence integration framework}
that assigns each candidate compound a calibrated posterior probability:

\[P(\text{compound}_i \mid D) \propto \pi_i \cdot \text{LR}_{\text{MS2}} \cdot \text{LR}_{\text{MS1}} \cdot \text{LR}_{\text{RT}}\]
xs An explicit novel compound hypothesis is always included so all
posteriors sum to 1:

\[\sum_i P(\text{compound}_i \mid D) + P(\text{novel} \mid D) = 1\]

When no candidate is convincingly better than the novel hypothesis, the
model says noThe pipeline computes InChIKey14 from SMILES for each
library hit, computes it for the confirmed compound from the spectra
file's SMILES, and marks a hit correct=True if they match. rather than
forcing a call.

\subsubsection{1.2 What This Is Not}\label{what-this-is-not}

\begin{itemize}
\tightlist
\item
  \textbf{Not a binary classifier.} There is no FP ground truth --- all
  confirmed spectra were annotated because they had good library
  matches. The model has no labeled negatives. YY/ZZ could serve as fp
  but they are labelled not because of their fp nature but being weird
  as a spectra.
\item
  \textbf{Not a re-ranking model.} Rank alone tells you nothing about
  the margin between candidates. The model's value is in quantifying
  uncertainty, not improving rank.
\item
  \textbf{Not a general identification engine.} The model operates only
  on spectra that return library hits. It does not address de novo
  structure elucidation.
\end{itemize}

\subsubsection{1.3 Why Bayesian}\label{why-bayesian}

The Bayesian framework was chosen not for accuracy but for properties
that matter scientifically:

\begin{enumerate}
\def\labelenumi{\arabic{enumi}.}
\tightlist
\item
  \textbf{Calibrated output} --- a posterior of 0.9 should mean the call
  is correct \textasciitilde90\% of the time
\item
  \textbf{Channel dropout} --- when RT/predicted RT is unavailable,
  LR\_RT = 1 and drops out silently without breaking the model
\item
  \textbf{Decomposable} --- each channel's contribution is explicit;
  ablations show exactly what is and is not working
\item
  \textbf{Modular prior} --- the prior can be swapped or improved
  independently of the likelihood terms
\item
  \textbf{Directly interpretable} --- P(threonine \textbar{} this
  spectrum) = 0.976 is a statement that anyone could comprehend.
\end{enumerate}

\begin{center}\rule{0.5\linewidth}{0.5pt}\end{center}

\subsection{Part II --- Data}\label{part-ii-data}

\subsubsection{2.1 Spectra}\label{spectra}

\textbf{\texttt{ttof+neg+hilic.csv}} --- 7,530 LC-MS/MS features from a
HILIC negative-ion TTOF (Premier) run processed through LC-BinBase.

Key fields used: \texttt{wiki\_id}, \texttt{rt} (observed retention
time, seconds), \texttt{precursor\_mz}, \texttt{entropy} (spectral
entropy H of the MS2 spectrum).

\subsubsection{2.2 Library Hits --- Two
Sources}\label{library-hits-two-sources}

All library hits were fetched from MassWiki, which aggregates public and
in-house spectral libraries.

\textbf{Reference library} (public: NIST23, GNPS, MassBank) - Broad
chemical coverage, big, variable spectral quality - RT channel: Kong 2.0
predicted RT (\texttt{predicted\_rt\_hilic}) - Mean entropy similarity
on confirmed correct hits: 0.798

\textbf{Annotation library} (in-house) - Spectra measured on the same
instrument (under the same conditions like collision energy?) - RT
channel: real measured RT offset (\texttt{anno\_delta\_rt} = observed RT
− library RT) - Mean entropy similarity on confirmed correct hits: 0.989
- Higher similarity and tighter RT

When an annotation hit exists, it typically dominates. The two sources
receive separate RT likelihood parameters (σ\_RT\_anno = 14.1 s vs
σ\_RT\_ref = 19.1 s) reflecting this quality difference.

\subsubsection{2.3 Biological Prior}\label{biological-prior}

Built from 17,415 confirmed TTOF HILIC negative records in
\texttt{hilic\_ttof\_all.csv}, queried from BinBase. Each compound's
prior weight is:

\[\pi_{\text{raw}} = \left(\sum_{\text{records}} \text{peak\_gaussian\_similarity} \times \text{peak\_purity}\right) \times \sqrt{N_{\text{studies}}}\]

Quality-weighting by peak shape and purity downweights poorly resolved
peaks; the square root of study count rewards compounds confirmed across
multiple independent studies.

\begin{center}\rule{0.5\linewidth}{0.5pt}\end{center}

\subsection{Part III --- The Scoring
Model}\label{part-iii-the-scoring-model}

\subsubsection{3.1 MS2 Likelihood Ratio}\label{ms2-likelihood-ratio}

Entropy similarity (0--1) is converted to a likelihood ratio via
softmax, weighted by spectral complexity:

\[\lambda(H) = 1 - e^{-\alpha H}\]

\[\text{LR}_{\text{MS2},i} = \left[\lambda \cdot \frac{e^{s_i}}{\sum_j e^{s_j}} + (1-\lambda) \cdot \frac{1}{N}\right] \times N\]

For H ≈ 0 (single-ion spectrum), λ ≈ 0 --- MS2 is uninformative, the
model falls back to the prior. For H \textgreater{} 0.5, λ ≈ 1 --- MS2
is fully trusted. The parameter α = 10.0 was estimated by maximum
likelihood on confirmed spectra; it consistently hits the upper bound.
Only 134 / 1,631 (8.2\%) of confirmed spectra have H \textless{} 0.5,
suggesting that the complexity weighting is only meaningfully active for
the 134 spectra below that threshold. Even at H = 0.3, MS2 is still 95\%
weighted.

\subsubsection{3.2 MS1 Likelihood Ratio}\label{ms1-likelihood-ratio}

Gaussian on Δppm with σ\_M = 4.21 ppm fitted from confirmed correct
hits. The denominator represents a random unrelated compound (null mass
distribution, σ = 50 ppm):

\[\text{LR}_{\text{MS1},i} = \frac{\mathcal{N}(\Delta\text{ppm}_i; 0, 4.21)}{\mathcal{N}(0; 0, 50)}\]

Note: TTOF is less mass-accurate than Orbitrap (σ\_M ≈ 4.2 ppm vs 1--2
ppm). MassWiki pre-filters candidates by precursor m/z, so most
candidates for a given spectrum already have small Δppm values --- the
MS1 channel is most valuable when a candidate has an unusually large
error (e.g.~the 14.6 ppm hit in Example 2, Section 5.2).

\subsubsection{3.3 RT Likelihood Ratio --- Two
Channels}\label{rt-likelihood-ratio-two-channels}

\textbf{Annotation hits} (real measured RT):
\[\text{LR}_{\text{RT},i} = \frac{\mathcal{N}(\texttt{anno\_delta\_rt}_i; 0, 14.1\,\text{s})}{\mathcal{N}(0; 0, 300\,\text{s})}\]

\textbf{Reference hits} (Kong predicted RT):
\[\text{LR}_{\text{RT},i} = \frac{\mathcal{N}(RT_{\text{obs}} - \widehat{RT}_i; 0, 19.1\,\text{s})}{\mathcal{N}(0; 0, 300\,\text{s})}\]

If no RT is available, LR\_RT = 1 (channel drops out). The 300 s null
represents a compound eluting at a random position on the gradient for a
5 min method (which is what we are using i suppose?).

\subsubsection{3.4 Novel Compound Term and
Abstention}\label{novel-compound-term-and-abstention}

\[P(\text{novel} \mid D) = \frac{P_{\text{novel}}}{\sum_i \text{unnorm}_i + P_{\text{novel}}}\]

If P(novel \textbar{} D) ≥ 0.6, the model abstains. P\_novel = 0.224 is
estimated from the fraction of all 7,530 MassWiki download spectra with
no library hit.

\subsubsection{3.5 Fitted Parameters}\label{fitted-parameters}

All parameters estimated from the 1,631 confirmed, human annotated
spectra:

{\def\LTcaptype{none} % do not increment counter
\begin{longtable}[]{@{}
  >{\raggedright\arraybackslash}p{(\linewidth - 4\tabcolsep) * \real{0.3333}}
  >{\raggedright\arraybackslash}p{(\linewidth - 4\tabcolsep) * \real{0.3333}}
  >{\raggedright\arraybackslash}p{(\linewidth - 4\tabcolsep) * \real{0.3333}}@{}}
\toprule\noalign{}
\begin{minipage}[b]{\linewidth}\raggedright
Parameter
\end{minipage} & \begin{minipage}[b]{\linewidth}\raggedright
Value
\end{minipage} & \begin{minipage}[b]{\linewidth}\raggedright
Method
\end{minipage} \\
\midrule\noalign{}
\endhead
\bottomrule\noalign{}
\endlastfoot
σ\_M & 4.21 ppm & Trimmed std of Δppm on InChIKey14-matched correct
hits \\
σ\_RT\_anno & 14.1 s & Std of \texttt{anno\_delta\_rt} on confirmed
annotation hits, trimmed ±60 s \\
σ\_RT\_ref & 19.1 s & Std of obs\_RT − predicted\_RT on confirmed
reference hits, trimmed ±120 s \\
α & 10.0 & MLE on confirmed spectra (hits optimizer upper bound) \\
P\_novel & 0.224 & Fraction of all 7,530 spectra with no library hit \\
\end{longtable}
}

\begin{center}\rule{0.5\linewidth}{0.5pt}\end{center}

\subsection{Part IV --- Performance on Confirmed
Spectra}\label{part-iv-performance-on-confirmed-spectra}

\subsubsection{4.1 Overall Accuracy}\label{overall-accuracy}

\textbf{Dataset:} 1,631 confirmed spectra (correct compound present in
library for all).

{\def\LTcaptype{none} % do not increment counter
\begin{longtable}[]{@{}ll@{}}
\toprule\noalign{}
Method & Top-1 Accuracy \\
\midrule\noalign{}
\endhead
\bottomrule\noalign{}
\endlastfoot
Raw entropy similarity rank & 99.8\% \\
Full Bayesian model & 90.3\% \\
LOO cross-validation & 90.7\% \\
Abstentions & 22 / 1,631 (1.3\%) \\
\end{longtable}
}

The 9.5 pp drop from the entropy baseline is expected and not a failure.
The confirmed set was annotated via library search --- entropy
similarity nearly always ranks the correct compound first by
construction. The RT and prior channels introduce noise on this set
because the prior is nearly flat across most candidates and σ\_RT
(14--19 s) is wide relative to RT differences between structural
isomers. The baseline's 99.8\% is a ceiling artifact, not a true
performance benchmark.

\subsubsection{4.2 Calibration}\label{calibration}

The model's primary contribution is not re-ranking but calibrated
uncertainty. When the model is confident, it is right:

{\def\LTcaptype{none} % do not increment counter
\begin{longtable}[]{@{}lll@{}}
\toprule\noalign{}
Posterior bin & N spectra & Observed accuracy \\
\midrule\noalign{}
\endhead
\bottomrule\noalign{}
\endlastfoot
\textgreater{} 0.90 & 93 & \textbf{100\%} \\
0.70--0.90 & 145 & \textbf{99.3\%} \\
0.50--0.70 & 253 & \textbf{97.2\%} \\
0.30--0.50 & 418 & \textbf{93.3\%} \\
\textless{} 0.30 & 700 & \textbf{83.9\%} \\
\end{longtable}
}

The monotonic relationship between posterior and accuracy is what makes
the score actionable. A threshold of 0.70 gives 99\%+ precision on this
dataset. How we generalize this to other datasets or other methods
remains unknown.

\subsubsection{4.3 The Duplicate Entry
Problem}\label{the-duplicate-entry-problem}

\textbf{87.9\% of confirmed spectra have more than one
\texttt{correct}-labeled library entry} (mean 6.0 per spectrum). The
same compound appears under multiple names, --- all sharing the same
InChIKey14. The model scores these as independent candidates, splitting
posterior mass.

Example: Threonine has 41 correct entries across the library. The top
posterior per entry is 0.063 --- which looks like low confidence --- but
the sum across all threonine entries is 0.98.

\textbf{InChIKey14 deduplication before scoring} collapses all entries
with the same molecular connectivity into one candidate before scoring
would eliminate posterior dilution. The pipeline computes InChIKey14
from SMILES for each library hit, computes it for the confirmed compound
from the spectra file's SMILES, and marks a hit \texttt{correct=True} if
they match.

\subsubsection{4.4 Failure Analysis}\label{failure-analysis}

159 of 1,631 confirmed spectra are called incorrectly. Of these: -
\textbf{68 (43\%):} correct compound is rank 2 --- near-tie, median
posterior gap = 0.084 - \textbf{91 (57\%):} correct compound is rank 3
or lower

Representative failure patterns:

\textbf{Name fragmentation} (\texttt{a6ICAHB/4704}): The correct
compound Fuc(a1-2)Gal(b1-4)GlcNAc has perfect entropy similarity (1.000)
but its probability is split with synonyms. Lewis A trisaccharide --- a
different name for essentially the same compound --- takes the top
posterior (0.775) while the correctly-named entry scores 0.096.

\textbf{Lipid notation ambiguity} (\texttt{a6ICAHB/4072}): PE
O-18:1\_16:1 and PE O-16:1\_18:1 are sn-positional isomers of the same
ether-linked phosphatidylethanolamine. The model calls PE O-18:1\_16:1
(wrong sn-position assignment) at 0.645. The correct entry PE
O-16:1\_18:1 scores 0.103. MS2 cannot distinguish sn-positions --- RT or
additional fragmentation would be needed.

\textbf{Synonym conflict} (\texttt{a6ICAHB/900}):
1,2-Dioleoyl-sn-glycero-3-phosphate (wrong) scores 0.576 while its
synonym 1-Oleoyl-L-α-lysophosphatidic acid (correct) scores 0.081 and
LPA 18:1 (correct) scores 0.013. All three refer to the same compound.

\begin{center}\rule{0.5\linewidth}{0.5pt}\end{center}

\subsection{Part V --- Expansion to Unconfirmed
Spectra}\label{part-v-expansion-to-unconfirmed-spectra}

\subsubsection{5.1 Motivation}\label{motivation}

The initial confirmed pipeline is inherently circular: fitted on
confirmed spectra, evaluated on confirmed spectra, which were themselves
selected because they had good library matches. The 5,899 unconfirmed
spectra had never been queried against MassWiki. Expanding the search:

\begin{enumerate}
\def\labelenumi{\arabic{enumi}.}
\tightlist
\item
  Provides a genuine test of the model on spectra selected independently
  of library search quality
\item
  Generates a prioritised annotation queue to replace arbitrary manual
  curation
\item
  Establishes the true P\_novel across the full run
\end{enumerate}

\subsubsection{5.2 Library Hit Coverage --- Correcting
P\_novel}\label{library-hit-coverage-correcting-p_novel}

{\def\LTcaptype{none} % do not increment counter
\begin{longtable}[]{@{}lll@{}}
\toprule\noalign{}
Category & N & \% of 7,530 total \\
\midrule\noalign{}
\endhead
\bottomrule\noalign{}
\endlastfoot
Confirmed spectra & 1,631 & 21.7\% \\
Unconfirmed with library hits & 4,216 & 56.0\% \\
True no-hit (no library match) & 1,683 & 22.4\% \\
\end{longtable}
}

\textbf{True P\_novel = 0.224.} The previous estimate of 0.783 was
entirely circular --- it counted unqueried spectra as ``no hit.'' Most
unconfirmed spectra do have library candidates; they had simply never
been scored.

\subsubsection{5.3 Scoring Results}\label{scoring-results}

167,857 candidate hits across 4,216 spectra were scored using the same
parameters fitted on the confirmed set (mean 39.8 candidates per
spectrum vs 17.5 for confirmed spectra --- the unconfirmed set has
broader and noisier library matches).

{\def\LTcaptype{none} % do not increment counter
\begin{longtable}[]{@{}lll@{}}
\toprule\noalign{}
Outcome & N & \% \\
\midrule\noalign{}
\endhead
\bottomrule\noalign{}
\endlastfoot
Called (P\_novel \textless{} 0.6) & 3,119 & 74.0\% \\
Abstained (P\_novel ≥ 0.6) & 1,097 & 26.0\% \\
\end{longtable}
}

Confidence distribution of called spectra:

{\def\LTcaptype{none} % do not increment counter
\begin{longtable}[]{@{}ll@{}}
\toprule\noalign{}
Posterior bin & N \\
\midrule\noalign{}
\endhead
\bottomrule\noalign{}
\endlastfoot
\textgreater{} 0.90 & 110 \\
0.70--0.90 & 82 \\
0.50--0.70 & 117 \\
0.30--0.50 & 327 \\
\textless{} 0.30 & 2,483 \\
\end{longtable}
}

The distribution is heavily skewed toward low confidence. The
unconfirmed set contains many genuinely ambiguous cases with no
annotation to guide selection. However, we identified 190 high
confidence calls that were not previous annotated for analytical
chemists to review.

\subsubsection{5.4 Review Queue}\label{review-queue}

2,253 spectra flagged across four categories for chemist review:

\begin{center}\rule{0.5\linewidth}{0.5pt}\end{center}

\paragraph{Flag 1 --- HIGH\_CONF\_NOVEL (156
spectra)}\label{flag-1-high_conf_novel-156-spectra}

\emph{Posterior \textgreater{} 0.70 on a compound not previously
confirmed in BinBase}

{\def\LTcaptype{none} % do not increment counter
\begin{longtable}[]{@{}
  >{\raggedright\arraybackslash}p{(\linewidth - 12\tabcolsep) * \real{0.1429}}
  >{\raggedright\arraybackslash}p{(\linewidth - 12\tabcolsep) * \real{0.1429}}
  >{\raggedright\arraybackslash}p{(\linewidth - 12\tabcolsep) * \real{0.1429}}
  >{\raggedright\arraybackslash}p{(\linewidth - 12\tabcolsep) * \real{0.1429}}
  >{\raggedright\arraybackslash}p{(\linewidth - 12\tabcolsep) * \real{0.1429}}
  >{\raggedright\arraybackslash}p{(\linewidth - 12\tabcolsep) * \real{0.1429}}
  >{\raggedright\arraybackslash}p{(\linewidth - 12\tabcolsep) * \real{0.1429}}@{}}
\toprule\noalign{}
\begin{minipage}[b]{\linewidth}\raggedright
wiki\_id
\end{minipage} & \begin{minipage}[b]{\linewidth}\raggedright
Compound
\end{minipage} & \begin{minipage}[b]{\linewidth}\raggedright
Posterior
\end{minipage} & \begin{minipage}[b]{\linewidth}\raggedright
H
\end{minipage} & \begin{minipage}[b]{\linewidth}\raggedright
Entropy sim
\end{minipage} & \begin{minipage}[b]{\linewidth}\raggedright
Δppm
\end{minipage} & \begin{minipage}[b]{\linewidth}\raggedright
Source
\end{minipage} \\
\midrule\noalign{}
\endhead
\bottomrule\noalign{}
\endlastfoot
a6ICAHB/6660 & Thr-Asp & 0.999 & 1.23 & 0.026 & −1.71 & annotation \\
a6ICAHB/4929 & D(+)-Trehalose & 0.999 & 2.13 & 0.263 & −0.41 &
reference \\
a6ICAHB/2469 & Atranone A & 0.999 & 2.17 & 0.055 & +1.87 & annotation \\
a6ICAHB/5767 & Glycodeoxycholic acid 3-sulfate & 0.999 & 0.14 & 0.344 &
−3.02 & annotation \\
a6ICAHB/243 & PC 18:0 & 0.998 & 1.21 & 0.327 & −2.88 & annotation \\
\end{longtable}
}

These are the highest-priority review candidates. If RT is chemically
plausible, each represents a new confirmed annotation. Note that several
entries have low entropy similarity (e.g.~Thr-Asp: 0.026) --- confidence
is driven by RT and MS1 alignment rather than spectral matching. These
should be verified with extra care.

\begin{center}\rule{0.5\linewidth}{0.5pt}\end{center}

\paragraph{Flag 2 --- NEAR\_TIE (1,188
spectra)}\label{flag-2-near_tie-1188-spectra}

\emph{Top two candidates within 0.05 posterior of each other, top
posterior \textgreater{} 0.10}

The model cannot resolve between two candidates. The most common cause
is structural isomers or lipid species with similar masses and RT
predictions. An authentic standard or manual RT inspection would resolve
most cases. This list is most useful for \textbf{prioritising standard
purchases} --- buy standards for the compounds that appear most
frequently in near-tie spectra.

\begin{center}\rule{0.5\linewidth}{0.5pt}\end{center}

\paragraph{Flag 3 --- RICH\_NO\_HIT (845
spectra)}\label{flag-3-rich_no_hit-845-spectra}

\emph{Spectral entropy H \textgreater{} 1.0, model abstains}

Well-fragmented spectra (median H = 1.65, max H = 3.55) with no
convincing library match. These are the most scientifically interesting
for discovery --- the spectrum contains real chemical information but
nothing in MassWiki explains it.

Example: \texttt{a6ICAHB/4373} has H = 3.547 (extremely
information-rich) but P\_novel = 0.999. The best library candidate (a
pyrimido-benzazepine derivative) has entropy similarity 0.178 and a 98 s
RT mismatch --- clearly wrong. This spectrum is a strong candidate for
de novo structure elucidation.

\begin{center}\rule{0.5\linewidth}{0.5pt}\end{center}

\paragraph{Flag 4 --- MS1\_RESCUE (64
spectra)}\label{flag-4-ms1_rescue-64-spectra}

\emph{Entropy similarity \textless{} 0.50 but posterior \textgreater{}
0.60}

RT and mass accuracy are doing the heavy lifting; spectral matching is
poor. These calls are the most suspicious and warrant explicit
verification.

Example: \texttt{a6ICAHB/4008} is called as γ-Glu-Leu at posterior 0.997
with entropy similarity of only 0.051. The near-zero mass error (−0.14
ppm) and annotation RT match drive the call. A 23 s ΔRT on an annotation
hit is within the fitted σ\_RT\_anno = 14.1 s window but is on the wide
side --- the RT trace should be inspected.

\begin{center}\rule{0.5\linewidth}{0.5pt}\end{center}

\subsection{Part VI --- Limitations and Known
Issues}\label{part-vi-limitations-and-known-issues}

\subsubsection{6.1 Confirmed Set Bias}\label{confirmed-set-bias}

All parameters (σ\_M, σ\_RT, α) are fitted on confirmed spectra, which
were confirmed because they had good library matches. This biases: -
\textbf{α} toward the spectral quality of curated spectra --- α always
hits its upper bound because confirmed spectra almost always have good
MS2 signal - The training set has no genuine FP ground truth - How can
we make sure the confirmed spectra are ``confirmed'' enough that no
errors are within this set?

\subsubsection{6.2 Annotation Library RT
Offset?}\label{annotation-library-rt-offset}

σ\_RT\_anno = 14.1 s is wider than expected for same-instrument
measurements, and the median annotation ΔRT is +9.2 s may suggest a
systematic shift? (This does not make sense to me as those rts are
corrected via istds). This may indicates a run-to-run RT drift or
gradient calibration difference between the annotation library and the
current run. The Gaussian model assumes zero-mean error; the actual
distribution is biased.

\subsubsection{6.3 Duplicate Library
Entries}\label{duplicate-library-entries}

As described in Section 4.3, 87.9\% of confirmed spectra have multiple
entries for the same compound. This dilutes posteriors and artificially
lowers per-entry confidence. InChIKey14 deduplication before scoring
would be helpful without any new data.

\subsubsection{6.4 α Is Unidentifiable}\label{ux3b1-is-unidentifiable}

α always hits the upper bound (10.0) in every LOO(leave one out) fold.
The confirmed spectra provide no leverage to estimate spectral
complexity weighting because nearly all confirmed spectra have H
\textgreater{} 0.5 (where λ ≈ 1 regardless of α). α is effectively
fixed, not learned.

\subsubsection{6.5 Initial P\_novel Was
Circular}\label{initial-p_novel-was-circular}

The original P\_novel = 0.783 was computed from the fraction of spectra
with no library hit, but only confirmed spectra had been queried.
Unconfirmed spectra were guaranteed to have no hits. The correct
P\_novel = 0.224 substantially changes the abstention behavior --- with
a lower novel term, weak library evidence is less easily suppressed.

\subsubsection{6.6 High-Confidence Wrong Calls in Unconfirmed
Set}\label{high-confidence-wrong-calls-in-unconfirmed-set}

Some HIGH\_CONF\_NOVEL calls have high posterior driven by RT/MS1 alone
with very low entropy similarity. Without a confirmed annotation or
authentic standard, these cannot be verified. Any automated acceptance
of these calls would require an independent check (e.g.~RT consistency
with the compound's chemical class on HILIC, or cross-referencing
against another study).

\begin{center}\rule{0.5\linewidth}{0.5pt}\end{center}

\subsection{Part VII --- Recommended
Use}\label{part-vii-recommended-use}

\subsubsection{Confirmed spectra}\label{confirmed-spectra}

{\def\LTcaptype{none} % do not increment counter
\begin{longtable}[]{@{}ll@{}}
\toprule\noalign{}
Posterior & Action \\
\midrule\noalign{}
\endhead
\bottomrule\noalign{}
\endlastfoot
\textgreater{} 0.90 & Accept identification (100\% precision on this
dataset) \\
0.70--0.90 & Accept with note for review \\
0.50--0.70 & Report tentative; include top 2--3 candidates \\
\textless{} 0.30 or abstain & Report as unidentified \\
\end{longtable}
}

\subsubsection{Unconfirmed spectra review
queue}\label{unconfirmed-spectra-review-queue}

{\def\LTcaptype{none} % do not increment counter
\begin{longtable}[]{@{}
  >{\raggedright\arraybackslash}p{(\linewidth - 6\tabcolsep) * \real{0.2500}}
  >{\raggedright\arraybackslash}p{(\linewidth - 6\tabcolsep) * \real{0.2500}}
  >{\raggedright\arraybackslash}p{(\linewidth - 6\tabcolsep) * \real{0.2500}}
  >{\raggedright\arraybackslash}p{(\linewidth - 6\tabcolsep) * \real{0.2500}}@{}}
\toprule\noalign{}
\begin{minipage}[b]{\linewidth}\raggedright
Priority
\end{minipage} & \begin{minipage}[b]{\linewidth}\raggedright
Flag
\end{minipage} & \begin{minipage}[b]{\linewidth}\raggedright
N
\end{minipage} & \begin{minipage}[b]{\linewidth}\raggedright
Action
\end{minipage} \\
\midrule\noalign{}
\endhead
\bottomrule\noalign{}
\endlastfoot
1 & HIGH\_CONF\_NOVEL, post \textgreater{} 0.90 & 110 & Verify RT; if
plausible, accept as new annotation \\
2 & MS1\_RESCUE & 64 & Inspect RT trace; low entropy sim calls need
verification \\
3 & RICH\_NO\_HIT, H \textgreater{} 2.0 & \textasciitilde200 & De novo
elucidation queue \\
4 & NEAR\_TIE & 1,188 & Prioritise authentic standard purchases \\
\end{longtable}
}

\begin{center}\rule{0.5\linewidth}{0.5pt}\end{center}

\subsection{Part VIII --- Proposed Next
Steps}\label{part-viii-proposed-next-steps}

{\def\LTcaptype{none} % do not increment counter
\begin{longtable}[]{@{}
  >{\raggedright\arraybackslash}p{(\linewidth - 2\tabcolsep) * \real{0.5000}}
  >{\raggedright\arraybackslash}p{(\linewidth - 2\tabcolsep) * \real{0.5000}}@{}}
\toprule\noalign{}
\begin{minipage}[b]{\linewidth}\raggedright
Improvement
\end{minipage} & \begin{minipage}[b]{\linewidth}\raggedright
Expected impact
\end{minipage} \\
\midrule\noalign{}
\endhead
\bottomrule\noalign{}
\endlastfoot
InChIKey14 deduplication before scoring & Eliminates posterior dilution;
likely pushes confirmed accuracy above 95\% \\
Correct/address annotation library RT offset (+9.2 s systematic shift) &
Fixes suppression of correct annotation hits; most impactful for
same-platform identifications \\
Fetch all 7,530 spectra (not just annotated) for parameter fitting?
maybe? still thinking about it & \\
Library source quality weighting/biological plausibility weighting (how
often is this chemical detected in the same matrix accross studies?
/chemical plausibility weighting (pubmed/patent count) & Better
priors \\
SIRIUS fingerprint as additional channel & Independent
spectral-to-structure prediction; genuinely orthogonal to entropy
similarity \\
Mixture model on full score distribution & Learn TP/FP distributions
from all spectra, not just confirmed ones; no FP ground truth needed \\
\end{longtable}
}

\begin{center}\rule{0.5\linewidth}{0.5pt}\end{center}

\subsection{Appendix --- Output Files}\label{appendix-output-files}

{\def\LTcaptype{none} % do not increment counter
\begin{longtable}[]{@{}
  >{\raggedright\arraybackslash}p{(\linewidth - 4\tabcolsep) * \real{0.3333}}
  >{\raggedright\arraybackslash}p{(\linewidth - 4\tabcolsep) * \real{0.3333}}
  >{\raggedright\arraybackslash}p{(\linewidth - 4\tabcolsep) * \real{0.3333}}@{}}
\toprule\noalign{}
\begin{minipage}[b]{\linewidth}\raggedright
File
\end{minipage} & \begin{minipage}[b]{\linewidth}\raggedright
Description
\end{minipage} & \begin{minipage}[b]{\linewidth}\raggedright
Rows
\end{minipage} \\
\midrule\noalign{}
\endhead
\bottomrule\noalign{}
\endlastfoot
\texttt{out\_proposal/assertions.csv} & Per-(spectrum, candidate)
posteriors, confirmed set & 28,510 \\
\texttt{out\_proposal/top\_calls.csv} & Top call per confirmed spectrum
& 1,631 \\
\texttt{out\_unconfirmed/unconfirmed\_scored.csv} & Per-(spectrum,
candidate) posteriors, unconfirmed set & 167,857 \\
\texttt{out\_unconfirmed/unconfirmed\_top\_calls.csv} & Top call per
unconfirmed spectrum & 4,216 \\
\texttt{out\_unconfirmed/review\_queue.csv} & Flagged spectra for
chemist review & 2,253 \\
\texttt{data/hilic\_ttof\_neg\_masswiki\_hits\_new.csv} & Library hits
for confirmed spectra & 120,048 \\
\texttt{data/hilic\_ttof\_neg\_masswiki\_hits\_unconfirmed.csv} &
Library hits for unconfirmed spectra & 167,857 \\
\end{longtable}
}
